% ============================================================
% ch5_conclusion.tex — 第五章 结论与展望
% 对齐 docs/thesis_chapter_drafts_20260420.md §6+§7
% 保留原 unnumbered chapter 结构（\chapter* + TOC handling）
% ============================================================

\chapter*{结论}
\addcontentsline{toc}{chapter}{结论}
\markboth{结论}{结论}
\setcounter{chapter}{5}
\setcounter{section}{0}
\label{chap:conclusion}

本章回扣第一章提出的三个问题。前四章已经完成问题界定、方法实例化与实证验证；本章归总这些结论的学术价值、适用边界和后续方向。

% ============================================================
\section{研究结论}
\label{sec:conclusion-answers}% backward-compat with earlier \ref
% ============================================================

本节不回放图表或模型个案，只概括本文对三个问题的回答。

第一，行为引导校准可以在 INT8 基准路径上形成稳定、可复现且可审计的基础保真路径。第四章的基准路径验证表明，在最保守的位宽和受控协议下，这一原则能够在不依赖额外在线自适应的前提下保持与全精度高度一致的任务行为；官方 LongBench 真实数据对照只承担评测协议一致性检查，说明主文合成评测在有限范围内与真实数据对照方向一致，但不能被读成已经覆盖更广泛真实应用分布。这里的结论不是“INT8 本身容易做好”，而是行为引导原则已经在最干净的验证实例上落成了一条可信的基础保真路径。

第二，行为引导原则不能机械地下推到对称 \texttt{INT4}，但可以在角色感知的非对称路径中恢复为一条可用且可审计的低比特实例。第四章已经表明，对称 \texttt{INT4} 面对的不是温和退化，而是具有明显架构依附性的阶跃失稳，其中 Key 精度下降是更主要的触发侧。\texttt{INT4-RoleAlign} 的意义也不在于宣称低比特问题已经普遍解决，而是在当前已验证的同格式条件和角色感知设定下，给出一条从对称失效走向可用恢复的清晰路径。因此，本文对低比特问题的回答不是“INT4 必然成功”，而是只有在承认 K/V 角色差异、并据此组织格式与离线校准纪律时，低比特路径才可能稳定建立起来。

第三，由校准链路产出的行为敏感度画像不仅能够服务参数选择，也可以作为跨模型预算分配的组织依据，并在当前已验证的预算范围内揭示一张由模型族、规模与任务共同决定的跨模型策略适用区间图谱。分配器与 \texttt{AutoK} 的实证结果没有导向一个脱离条件的单一赢家；更稳定的事实是，不同模型会落入不同的高性能区间：有的更依赖早层保护，有的呈现中等规模的共识解，也有的表现为多个高性能簇并存而没有稳定最优。heuristic 作为强基线并不是与本文主线相冲突的例外；它说明在某些区间里，可行层集本身就很窄。行为引导方法的价值，在于识别这些区间何时出现、如何区分，并把原本分散的跨模型现象整理成可命名、可解释的结构读数，而不是把单一策略写成普适法则。

这三个回答共同指向一个更高层的判断：注意力行为可以作为贯通校准与分配的组织对象。它先在 INT8 基准路径上建立基础保真路径，随后在低比特场景下具体化为一条承认 K/V 角色差异的恢复路径，并进一步扩展成能够读出跨模型策略适用区间的结构化工具。第五章后续不再重复第四章的证据，而转向说明这些研究结论的创新点、学术价值与适用边界。

% ============================================================
\section{创新点与学术价值}
\label{sec:ch5-innovation-value}
% ============================================================

本节说明这些结论的学术价值。

理论层面的价值在于，本文没有把 KV Cache 量化继续拆成几个彼此独立的局部优化问题，而是把注意力行为作为同时组织校准与分配的共同对象。第三章从注意力近似误差传播分析出发，说明量化损伤被下游感知时，已经不再是孤立的张量数值偏差，而是由注意力分布偏移与聚合输出偏移共同构成的行为失真。校准层与分配层因此不再只是两个并列的工程步骤，而被收束到同一条行为引导主线之下。本文不是为既有量化流程再补一个局部校准目标，而是为 KV Cache 量化提供了一种能够同时容纳基准验证、低比特实例和预算扩展的解释框架。

方法层面的价值在于，同一原则下形成了一条层次分明的实例链。作为最保守也最干净的验证实例，INT8 基准路径说明行为引导校准能够在基准位宽上建立稳定、可复现的基础保真路径；在与 \texttt{KIVI-style} 共享非对称格式的前提下，\texttt{INT4-RoleAlign} 的意义也不在经验上的一边倒优势，而在于它把低比特路径组织成了一条可固化、可审计、并能继续向后扩展的实例化路径。分配器与 \texttt{AutoK} 不是脱离校准层另起炉灶的第二套方法，而是由同行为敏感度画像延伸出的预算决策扩展。本文把基准验证、低比特恢复与预算扩展组织成了一条前后衔接的分层实例链。

经验层面的价值在于识别并命名了一张跨模型策略适用区间图谱，而不是把某一种策略写成跨模型、跨任务、跨预算条件下的单一最优。第四章的跨模型结果并没有把比较收束成一个脱离条件的普适赢家；更稳定的事实是，模型家族、参数规模与任务会共同塑造不同策略的高性能区间。本文提供的也不是某个分配器在若干读数上赢了多少，而是一种用来辨认这些区间何时出现、如何分化以及为何分化的结构化视角。heuristic 在多个区间中表现为强基线，也并不削弱本文主线；它提醒我们，在某些模型条件下，可行层集本来就很窄，而行为引导方法的学术价值，正在于把“强基线何时成立、何时失效”这类现象从经验杂讯提升为可解释的结构读数。

因此，本文的创新点与学术价值，并不在于孤立地提出某一个新的量化技巧，而在于以注意力行为为统一组织对象，把基准路径、低比特实例与跨模型策略适用区间图谱组织成一条具有解释力、可审计性与可扩展性的研究路径。也正是在这一意义上，本文不仅回答了 KV Cache 量化“如何压缩”的问题，也把这一问题从局部技巧比较推进到了组织原则与结构解释的层面。

% ============================================================
\section{研究局限性}
\label{sec:conclusion-limitations}
% ============================================================

本节界定上述价值成立的范围。它们不能脱离评测协议、比较口径、系统环境和时间外推条件而无条件成立；讨论局限性，也不是撤回前面的结论，而是把这些结论的适用域交代清楚。

从评测协议上看，本文的主要实证证据建立在统一协议下的 LongBench 风格合成基准任务上，官方 LongBench 真实数据对照仅在 INT8 基准路径上承担评测协议一致性的补充检验。这说明本文已经能够在协议内给出稳定证据，也通过一块受控的真实数据对照提供了有限范围内的一致性证据，但还不能把相关结论外推到更广泛的真实应用场景。尤其当输入分布、任务结构、提示风格与真实部署流程发生变化时，不同量化路径的相对表现及其适用边界仍可能随之改变。与此同时，本文虽然并列使用了 PPL、Needle 与任务级指标以降低单一度量带来的偏差，但这并不意味着已经穷尽了长上下文质量的全部观察维度。

在比较口径上，本文关于逐层预算分配与跨模型策略适用区间图谱的主要结论，建立在同量级 \texttt{INT4} 预算带内的经验可比性上，而不是严格的同预算比较。现有证据足以支持不同模型存在不同结构足迹、模型族/规模/任务共同塑造高性能区间，以及 heuristic 作为强基线所具有的结构性含义；但这些证据还不足以形式化证明某一单一策略在严格同预算条件下具有普适更优性，也不足以完成分配器相对 \texttt{KIVI-style} 的正式胜负裁决。因此，本文在分配器线上的主要贡献，应被理解为对跨模型策略适用区间的识别与组织，而不是对严格预算可比性问题的最终完成。

此外，本文不少结论本来就是条件性的，其解释力更多体现为在特定模型族、参数规模与任务类型下识别稳定结构，而不是给出脱离条件的普适规律。尤其在低比特部分，\texttt{INT4-RoleAlign} 与 \texttt{KIVI-style} 的比较虽然已被限制在同格式条件下，从而能够较干净地呈现校准理念的差异，但格式贡献、校准贡献与架构因素之间的作用仍未被完全正交解耦。因此，当前证据更支持将 \texttt{INT4-RoleAlign} 理解为一条在既定格式与已验证设定下成立的低比特恢复路径，而不是一条已经对所有模型族、规模与任务条件给出统一因果归因的普适方案。类似地，跨模型策略适用区间图谱揭示的也是条件性结构，而不是一个脱离模型与任务语境的普适规律。

系统层与动态决策层同样有明确的适用域限制。部署效率、性能交叉边界与显存收益的读数严格绑定于当前的 \texttt{NVIDIA H20}、固定软件栈、当前 backend 组合以及统一解码设定，因此不能被直接推广为跨硬件、跨 runtime 的一般定律。与此同时，\texttt{AutoK} 当前主要仍是模型级的静态 proposer，prompt 级自适应路由只呈现出弱改善或混合信号，尚不足以被写成正式结论；而结果溯源链保证的是结果的可追溯与可审计，而不是对未来 Transformer 变体或新长上下文注意力架构的自动泛化。换句话说，本文给出的，是一个在当前架构、协议和系统环境下可信、可审计的工作版本，而不是可以不经复核就直接迁移到未来环境的答案。

因此，本节讨论的局限性并不意味着前述结论失效，而是把它们成立的条件正式限定下来。当前关于基准路径保真、低比特恢复与跨模型策略适用区间图谱的结论，已经在评测协议、比较口径、条件适用域以及系统与时间外推边界上被清楚界定。下一节将在此基础上，只沿其中可直接转写为研究路线的三类方法性边界，进一步展开未来工作展望。

% ============================================================
\section{未来工作展望}
\label{sec:conclusion-future}
% ============================================================

上述边界也限定了后续工作的方向。后续最值得推进的，并不是另起新的问题域，而是继续沿 prompt 级策略选择、角色感知预算机制和严格预算可比性三个方向扩展行为引导框架的适用域。

一条值得优先检验的路线，是将当前以模型级为主的预算建议机制进一步推进到 prompt 级的条件化策略选择。现阶段的 \texttt{AutoK} 主要仍根据离线行为敏感度画像给出模型级建议，而 prompt 级自适应方案的探索只呈现出弱改善或混合信号，尚不足以被写成本文的正式结论。这意味着，现有跨模型策略适用区间图谱虽然已经能够为固定预算提供结构化读数，但它对输入层差异的响应仍然比较粗粒度。后续工作可以继续检验：如果把输入提示的结构表征与当前图谱中的高性能区间联系起来，是否能够形成稳定、可迁移的 prompt 级路由机制，使策略调用从“按模型给建议”进一步推进到“按输入做条件化选择”。

沿着同一主线往前走，另一个重要方向，是把当前只具接口形态的角色感知预算扩展推进为更成熟的角色感知分配器。本文已经表明，K/V 两侧在低比特下并不对称，角色感知的预算语言也因此具备明确的结构基础；但从现有证据看，这种扩展更多体现为“可表达、可试验、局部可受益”，而不是已经形成稳定跨模型占优的成熟机制。后续工作不应只是再提出新的分配名称，而应进一步把 Key 与 Value 两侧的预算描述、耦合关系和画像读数结合起来，使预算决策从当前较粗的逐层配置，推进为更细致、也更一致的联合分配机制。只有当这种角色差异进入分配器的核心决策中，行为引导框架在分配层的解释力与可复用性才可能进一步加强。

在比较口径层面，分配器线后续最需要补全的，也不是更多零散个案，而是更严格的同预算比较。第四章已经证明，在当前同量级 INT4 预算带下，不同模型确实会落入不同的策略适用区间，这一点支撑了图谱层面的结构判断；但这并不等于某一单一策略已经在严格同预算条件下完成了形式化胜负裁决。因此，未来工作的重点应放在设计更干净的严格预算可比性协议，以补全分配器相对现有基线在正式比较意义上的证据链。比较口径收紧之后，本文关于策略适用区间的解释力，才能继续向更可比较的结构主张推进，而不只是停留在经验上同量级预算的读数组织。

总之，本文后续最值得推进的，并不是平行扩张更多方法分支，而是沿同一行为引导主线继续向三个层面深化：在策略调用层面优先检验 prompt 级路由，在预算机制层面推进更成熟的角色感知分配器，在比较口径层面补全严格同预算比较。未来工作的重点因此也不在于寻找一个脱离条件的“全局最优”策略，而在于持续扩大这一框架的可解释性、可比较性与适用域。

% ============================================================
\section{结语}
\label{sec:conclusion-closing}
% ============================================================

面向长上下文高效推理，KV Cache 量化并不是单纯寻找更低 bit-width 的局部优化问题，而是一个需要统一组织原则的系统设计问题。本文以注意力行为为统一组织对象，将校准与分配两层决策纳入同一行为引导框架，并在基准实例、低比特路径与预算扩展之间建立了可相互印证的分层实例链。

本文没有把跨模型实验压成一个脱离条件的“全局最优”策略，而是识别并命名了跨模型策略适用区间图谱，使不同模型族、规模与任务条件下的结构差异能够被读出、被解释，也为后续研究留下了一种可以继续承接的结构读法。在此基础上，本文对这些结论的协议边界、比较口径与系统适用域作出了明确限定。论文的最终价值不在于提出一个永远最优的局部量化技巧，而在于为面向长上下文高效推理的 KV Cache 量化提供了一条以注意力行为为中心、可审计、可扩展并具备结构解释力的研究路径。
